\section{One coefficient field and auxiliary primes on the marked frame}

Let $L$ be the CM field of the given Galois link, $L^+$ its maximal
real subfield, and $G_E=\operatorname{Gal}(\Omega/E)$. Keep the
original link $k$, its marked frame $F$, its modulus $M\ne0$, and its
coefficient ring $O\subset\overline{\mathbb Q}_\ell$. Assume the
original $\mathrm{LinkHyp}$ and $F.\mathrm{Ok}(C_0)$,
$F.\mathrm{Big}(C_0)$. Write $\Psi_F$ for the residual character
on the frame and $p_0:G_L\twoheadrightarrow F.G_0$ for the given
map to its index-two subgroup. We construct the two realizations
used in A68M5. The coefficient field is chosen before the auxiliary
places, so the two runs can change their places over one residual
frame.

% BEGIN C69C1
\paragraph{C69C1. A common coefficient ring for two continuous twists.}
Let $a_2,a_3:G_L\to\overline{\mathbb Q}_\ell^{\times}$ be continuous.
There is a single finite extension $E/\mathbb Q_\ell$, inside this
algebraic closure, such that $O_1=\mathcal O_E$ contains $O$, all
values and inverse values of both twists, and the images of
$\mathcal O_L$ under every embedding $L\to\overline{\mathbb Q}_\ell$.

Here one must put the whole image of each character in a finite
field extension. The algebraic closure is covered by a countable
family of finite extensions. For example embed $\overline{\mathbb Q}$
in it. This countable subfield is dense: approximate the minimal
polynomial of any element by rational polynomials and use continuity
of the simple roots (equivalently Hensel's lemma in the local field
generated by the element). A sufficiently close root generates a
field containing the given element by Krasner's lemma. Finite
composita of fields of these roots give the countable family. Each
finite extension is complete and hence closed in the algebraic
closure. For a continuous character, the inverse images of the
multiplicative groups of these fields form a countable closed
subgroup cover of the compact Hausdorff group $G_L$. The Baire
theorem gives a subgroup with interior. It is open and has finite
index. Adjoin the values at finitely many coset representatives to
obtain a finite extension containing the entire image. The norm of
the image is a compact subgroup of the positive reals, hence trivial
by taking positive and negative powers. All character values are units.

Take the compositum for the two twists, the field generated by finite
$\mathbb Z_\ell$-module generators of $O$, and the fields generated
by the images of a $\mathbb Q$-basis of $L$ under its finitely many
embeddings. The generators of $O$ are integral over $\mathbb Z_\ell$,
as are the images of algebraic integers, so their respective rings
lie in $O_1$. The resulting $O_1$ is a complete DVR finite over
$\mathbb Z_\ell$ with finite residue field of characteristic $\ell$.
The inclusion $O\hookrightarrow O_1$ is integral, and the maximal
ideal of $O_1$ contracts to the maximal ideal of $O$. Thus it induces
the required field map $f:k_O\to k_{O_1}$, taking the residue of an
element to the residue of that very element. This proves both the
largeness and the residue compatibility in $\mathrm{CoefEx}$.
More generally every coefficient DVR in the source has
$\mathrm{IsCoef}$: it is finite torsion free over $\mathbb Z_\ell$,
complete in the $\ell$-adic topology, and $\ell$ is a nonunit by
integrality. Its maximal-ideal topology agrees with the $\ell$-adic
topology by the DVR factorization of $\ell$.
% END C69C1

% BEGIN C69C2
\paragraph{C69C2. The twist on the linked CM field.}
We prove the source's $\mathrm{KaOf}$ from the algebraicity of
cyclotomic powers ($\mathrm{AlgCyc}$) and preservation of algebraic
characters by restriction ($\mathrm{AlgRes}$). In rank three take
$a_3=\epsilon_\ell^{-1}|_{G_L}$. It is algebraic by the first input.
Complex conjugation fixes the cyclotomic character, so the product
of the two conjugate values is $\epsilon_\ell^{-2}$. It is unramified
away from $\ell$, hence outside $M$ since $\ell\mid M$, and has
inertia exponent one away from $\ell$.

For rank two suppose $G_K=\ker\eta$ inside $G_{\mathbb Q}$ and
$\kappa_K$ is the supplied $\mathrm{IsTwist}$ on $G_K$. The link's
fixed-group identity gives $G_L\subset G_K$, hence an inclusion
$K\hookrightarrow L$ compatible with the specified embeddings in
$\Omega$. The field $K$ is Galois over $\mathbb Q$ since $L/\mathbb Q$
is abelian. Every rational prime dividing $|\operatorname{disc}K|$
divides $M$. Indeed otherwise all its inertia groups lie in
$\ker\eta$ by $\mathrm{LinkHyp.ram}$, so their restrictions to $K$
are trivial, contradicting the discriminant criterion for ramification.
For a valuation ring of $\Omega$, each integer has valuation at most
one in multiplicative notation. An integer with valuation below one
has a prime divisor of valuation below one, by factorization. Two
coprime integers cannot both have valuation below one, by B\'ezout.
Thus outside $M$ one is simultaneously away from $\ell$ and
$\operatorname{disc}K$, including for the valuation-ring tests used
by the source.

Take $a_2=\kappa_K|_{G_L}$. Its algebraicity follows from
$\mathrm{AlgRes}$ and the algebraic-character data of $\mathrm{IsTwist}$.
Its conjugation identity is that of $\kappa_K$, with product
$\epsilon_\ell^{-1}$. Restriction sends $L$-inertia into $K$-inertia,
so unramifiedness outside $M$ follows from the preceding discriminant
argument. The single prime-to-$\ell$ exponent for all inertia away
from $\ell$ is the one supplied for $\kappa_K$. Finally the two
characters agree whenever the automorphisms have the same image in
$G_{\mathbb Q}$, by the compatible embedding of $K$ in $L$. These
are the full two $\mathrm{IsKa}$ conclusions with the required
rank-two restriction identity. The choice of $\kappa_K$ belongs to
the separate quadratic twist theorem.
% END C69C2

% BEGIN C69C3
\paragraph{C69C3. A degree-one place satisfying all the exclusions.}
Fix $t\ge1$. We choose the distinguished place before choosing the
level. On $F.G_0$ there is $x$ with $\Psi_F(x)=1$ and
$\bar\epsilon_\ell(x)\ne1$. Otherwise $\bar\epsilon_\ell$ factors
through $\Psi_F(F.G_0)$, a cyclic subgroup of $k_O^\times$. Every
character from that subgroup into $k_O^\times$ is a power of its
inclusion: the image of a generator has order dividing its order,
and the subgroup of that order is unique. This contradicts the
frame's condition f5. Lift $x$ to $g\in G_L$ by $p_0$.

In $G_{L^+}$ intersect three open normal subgroups: the kernel of
the pullback of $\Psi_F$ through the frame projection, the kernel
of the mod-$\ell$ cyclotomic character, and $G_L$. Apply Chebotarev
to the class of $g$ in this finite quotient and exclude all primes
over $2Mt|\operatorname{disc}L|$. One can impose residue degree one
over $\mathbb Q$: primes of degree greater than one contribute
$O([L^+:\mathbb Q]\sqrt X)$ primes of norm $\le X$, whereas the
specified Chebotarev class has positive density. Choose the place
above the resulting $v_0$ in $\Omega$ so its Frobenius maps to the
actual element represented by $g$ (conjugate the chosen place
within the finite Galois extension). Inertia lies in the threefold
intersection. Consequently $\operatorname{N}v_0=p$ is an odd rational
prime with $p\nmid\ell Mt\operatorname{disc}L$. The decomposition
group at $v_0$ is contained in $G_L$ and acts trivially through
$\Psi_F$, since inertia and Frobenius topologically generate it.
The same Frobenius has nontrivial mod-$\ell$ cyclotomic value. Thus
$v_0$ splits in $L/L^+$, the whole $L$-decomposition group maps
into $\ker\Psi_F$, and it contains the required cyclotomic witness.

Use the places over $\ell$, the finite set $S_t$ of places dividing
$M$ away from $\ell$, and $v_0$ as the level set $S$. Here $S_t$
denotes the source's tame set and is independent of the exclusion
integer $t$. The first two parts split by the link's splitting
condition at $M$. The link identifies its local indices with the
places of $L^+$ over $\ell$: prolongations above one such place
form a single $G_{L^+}$-orbit, and its places $k.\mathrm{pl}(v)$
are exactly one representative per such orbit by the link's
existence and uniqueness. Use those representatives at $\ell$,
the chosen prolongation at $v_0$, and any prolongations elsewhere.
At each place use the induced place on $L$. Take the complex
conjugation given by the link. Its unramifiedness condition says
that every inertia group over $L^+$ fixes $L$. These choices give
the Hida setting $s$, its $p$ and $v_0$, and all the linked place
identities. The cyclotomic and Dirichlet factors of the original
character are unramified outside $\ell N$, hence outside $M$.
KX2 gives a level $V$ with $V.S'=S$ exactly.
% END C69C3

% BEGIN C69C4
\paragraph{C69C4. The quotient, its local maps, and its kernel test.}
Fix a twist $a$ of either rank satisfying $\mathrm{IsKa}$ whose
values belong to $O_1$, and the level of C69C3. Put $H=G_{L^+}$,
$H^-=G_L$, and let $N_S$ be the closed normal subgroup of $H$
generated by inertia at places outside $S$. The level group is
$H/N_S$ and its sign kernel $G_0$ is the image of $H^-$. Write
$q:H^-\twoheadrightarrow G_0$ for that projection. All these
quotients carry their profinite quotient topologies. The global link
identifies $H$ with its restriction image in $G_{\mathbb Q}$.
There the kernel of its map to $F.\Gamma$ is generated as a closed
normal subgroup by the inertia elements away from $\ell$ which are
in $H^-$ and on which the residual arithmetic character is one.
Outside $S$ inertia fixes $L$ and this character is one by C69C3.
Thus $N_S$ is killed. Restricting the descended map gives
$\pi:G_0\twoheadrightarrow F.G_0$ with $\pi q=p_0$. Its
surjectivity follows from that of $p_0$. Conjugation and the
cyclotomic character agree since they agree upstairs and the chosen
complex conjugation is the link's own lift.

The character $a$ also descends to $\kappa:G_0\to O_1^\times$.
To check its kernel contains $N_S$, use all prolongations of all
places outside $S$. Each inertia group lies in $H^-$ and is killed
by $a$ by unramifiedness outside $M$. Their $H$-conjugates are
inertia groups of other prolongations above the same places and
hence are killed too. The kernel of $a$ is closed in $H^-$,
which is closed in $H$, so it contains their closed generated
subgroup. The quotient map $q$ transfers continuity to $\kappa$;
the norm topology on $O_1$ is its adic topology. In particular
$\kappa q=a$, and passage through the surjection gives
$\kappa(g)\kappa(cgc^{-1})=\epsilon_\ell(g)^{1-n}$.

At $v\mid\ell$ put $w=e(v)$, using C69C3's indexing. The arithmetic
decomposition group of the chosen prolongation and $F.G_w$ embed in
$G_{\mathbb Q}$ with the same image. Indeed the link identifies the
second image with the stabilizer in $H$ of $k.\mathrm{pl}(w)$, and
this entire stabilizer is in $H^-$ by splitting at $M$. Thus their
identification $\lambda_v$ is a topological group isomorphism by
compactness and Hausdorffness. On this common subgroup the local
maps to the frame agree by the link, as do the inertia subgroups.
In particular $k.j_w(\lambda_v u)$ represents the given $u$, and
the chosen valuation ring is exactly $k.\mathrm{pl}(w)$.

For the kernel test let $D\subset G_0$ be closed normal and stable
under conjugation by $c$. Its image is normal in $H/N_S$ since the
two cosets are represented by $1,c$. Its preimage in $H$ is closed
normal. Consider an inertia generator of the frame kernel described
above. At a place outside $S$ the generator is trivial in $H/N_S$.
At a place in $S_t$, conjugate its prolongation to the selected one.
The preimage contains it by the assumed tame condition on $D$, since
the extended residual character of an element with original residual
value one is one. At $v_0$ the assumed distinguished inertia inclusion
suffices. These cases exhaust generators away from $\ell$ (the trivial
valuation contributes trivial inertia). Normality and closedness put
the entire frame kernel in that preimage. Lifting through $q$ now
gives $\ker\pi\subset D$, exactly the link's universal kernel test.
% END C69C4

% BEGIN C69C5
\paragraph{C69C5. The residual determinant and the inertia conditions.}
Put $k'=k_{O_1}$, $\Psi'=f\Psi_F$, and let $w_i$ be the weights of
the rank-$n$ normalized frame. On the level group $G_0$ take
\[
 \bar\chi_i(g)=\Psi'(\pi g)^{w_i}\bar\kappa(g),\qquad
 \bar\psi_{vi}=\bar\chi_i\circ\phi_v,\qquad \theta_{vi}=1 ,
\]
where $\phi_v$ is the local projection. The kernels of the two
factors in $\bar\chi_i$ are open. Evaluation at a lift of the marked
element $F.g_0$ distinguishes these characters: $\Psi_F(F.g_0)^2
\ne1$, and the rank-two weights are $1,-1$ while the rank-three
weights are $1,0,-1$ in the frame order.
Under the weighted involution $[g]\mapsto
\bar\epsilon_\ell(g)^{1-n}[(cgc^{-1})^{-1}]$, evaluation by
$\bar\chi_i$ becomes evaluation by $\bar\chi_{n-1-i}$, since
$w_{n-1-i}=-w_i$ and C69C4 gives the product of the conjugate twists.
The same equality holds by linearity on every scalar extension of
the group algebra. Multiplying the diagonal evaluations proves
polarization as an identity of determinant laws.

At $u\in S_t$ use the single exponent $(|k_O|-1)m_a$, where $m_a$
is the prime-to-$\ell$ inertia exponent of $a$. Both factors are
prime to $\ell$. For an inertia lift $h$, its $(|k_O|-1)$-th
power is an inertia element with trivial residual character.
The closed normal kernel prescription in C69C4 kills this power
under $p_0$, while $a(h)^{m_a}=1$. This proves both tame power
conditions. At $v_0$ the twist is unramified since $p\nmid M$.
The entire decomposition group has $\Psi'=1$, so all
$\bar\chi_i$ agree there, each equals one on inertia, and the
witness of C69C3 still has nontrivial cyclotomic residue in $k'$,
as $\mathbb F_\ell\hookrightarrow k'$ is injective. These are the
three clauses of $r.\mathrm{V0ok}$.

At each $v\mid\ell$ transport the link's cyclotomic witness
through $\lambda_v$. Its nontrivial mod-$\ell$ action implies
the local field $L_w$ does not contain $\mu_\ell$, since that
action is on the actual $\ell$-th roots of unity. By local
reciprocity the topological abelianization of the decomposition
group is $\mathcal O_{L_w}^{\times}\times\widehat{\mathbb Z}$
after a choice of uniformizer (the local coordinates of KX3).
Its $\ell$-power torsion consists of the $\ell$-power roots of
unity in $L_w$, hence is trivial. Its inverse image $\mu_v$
therefore lies in the closed commutator subgroup, as required.

Finally the diagonal representation of the $\bar\chi_i$ gives a
residual ordinary pair. In each diagonal coordinate $j$ the ordered
product of local factors $\bar r(\phi_v(g_i))-\bar\psi_{vi}(g_i)I$
is zero because its $j$-th factor vanishes in that coordinate.
All local central characters on $\mu_v$ are trivial by the preceding
abelianization argument. At tame inertia in the common residual
kernel the matrix is the identity, so the required traces of its
difference with the identity vanish. Together with local constancy
and the polarized law above this gives a pair over $k'$ with exactly
the displayed global determinant and local characters. KX8 applies
to this pair and supplies the universal ordinary determinant ring
$B_0$ of this level and residual datum.
% END C69C5

% BEGIN C69C6
\paragraph{C69C6. The normalized matrix space over the enlarged ring.}
Set $F'=F.\mathrm{ext}(f)$. Its marked groups, local maps, sign,
global marking and cyclotomic character are those of $F$. Its
residual character is $\Psi'$. Injectivity of $f$ reflects equality
of characters, their powers and their kernels, transporting the
separation, nonfactorization and inertia conditions. For the
cohomology clauses, choose a $k_O$-basis of $k'$. A continuous
cochain with values in a transported one-dimensional character
module splits in that basis into finitely many continuous cochains
over $k_O$. Differentials and local restrictions act coordinatewise.
Thus the global and strict first cohomology groups extend scalars,
preserving their respective dimensions over the new field and their
nonvanishing. The Euler-characteristic clause already quantifies
over coefficient fields. All numerical inequalities and the
coefficient-independent hypotheses are inherited. This gives
$F'.\mathrm{Ok}(C_0)$ and $F'.\mathrm{Big}(C_0)$.

The matrix problem on $F'$ uses its untwisted diagonal residual
representation $\bar r=\mathrm{diag}(\Psi'^{w_i})$. Its entries are
locally constant. The finiteness condition on open subgroups
(Mazur's $\Phi_\ell$) holds on the open sign kernel of the frame
by restriction. Apply FP2 to get its universal unrestricted
continuous matrix lift $u$ over $U$. At each local index choose a
separator $s_v$ with $\Psi'(\mathrm{loc}_v s_v)^2\ne1$, from f3.
The matrix $u(\mathrm{loc}_v s_v)$ has simple labeled eigenvalues
modulo the maximal ideal. By Hensel lifting its characteristic
polynomial splits into linear factors with unit differences.
The spectral projectors onto these roots lift the standard
coordinate lines. Applying each projector to the corresponding
standard basis vector gives a matrix $P_v\equiv I$ which
diagonalizes the separator. Write
$T_v(h)=P_v^{-1}u(\mathrm{loc}_v h)P_v$.

Take the quotient of $U$ by the entries of
$u(g^*)-Ju(g)^tJ$ for all $g$, the off-diagonal entries of
$u(F.g_0)-u((F.g_0)^*)$, and the entries below the diagonal in
every $T_v(h)$. Here $J$ is the anti-diagonal identity and $g^*$
is the frame involution on group elements. The residual point kills
these entries, hence the quotient is a nonzero complete Noetherian
local ring. Its local diagonals give continuous characters. The
ordered product of the local factors for $h_0,\ldots,h_{n-1}$
vanishes: acting from the right, the last factor sends the full
standard flag space to its penultimate member, and each subsequent
factor drops the appropriate flag member one further step. This
gives a normalized lift with its ordered characters.

Conversely let a normalized lift with local characters $\psi_i$
give an unrestricted map out of $U$, and fix $v$. In the basis
transported from $P_v$ write $d_j$ for the diagonal entries of
the separator. In the ordered product put $h$ in position $i$
and $s_v$ in every other position. For $j\le i$ its $(i,j)$
entry factors as
\[
 \prod_{a<i}(d_i-\psi_a(s_v))\,
 [T_v(h)-\psi_i(h)I]_{ij}\,
 \prod_{a>i}(d_j-\psi_a(s_v))=0
\]
in the target. All outer factors are units by residual separation.
This forces triangularity and recovers each diagonal character.
The star and marked diagonal conditions kill the remaining
generators. The resulting map out of the quotient is unique by
FP2, since the characters have already been recovered. This is the
universal normalized matrix space $\mathcal S$ on $F'$ in the
source's sense.
% END C69C6

% BEGIN C69C7
\paragraph{C69C7. An arithmetic realization with a prescribed avoidance integer.}
For either rank use the concrete level and residual datum of
C69C3--C69C5, the twisted link of C69C4--C69C5, its determinant
ring $B_0$, and the space $\mathcal S$ of C69C6. The equations
$\pi q=p_0$ and $\kappa q=a$ identify the global projection and
twist on the original arithmetic group. At $\ell$ the maps
$\lambda_v$ identify the original embeddings and the actual
prolongations. The distinguished prime is away from $M$, has
$r.\mathrm{V0ok}$, and has rational norm $p\nmid t$.
Together with $O\subseteq O_1$, the residue compatibility of
C69C1 and the frame bounds, these are the fields of the original
$\mathrm{CReal}(k,O_1,f,n,a)$ on $F'$. For every nonzero natural
integer $t$ this proves the strengthened $\mathrm{CRealAv}$ and
in particular $\mathrm{CRealEx}$ by $t=1$.

For two runs of A68M5 take $t=1$ in both ranks, obtaining
$p_2,p_3$. On the second run take $t=p_2p_3$ in each rank. Each
new prime differs from each of the first two choices. Both runs
use the same inclusion, residue map, link, two arithmetic twists,
and marked frame. C69C1 takes continuous twists as inputs and
C69C2 uses the supplied algebraicity and quadratic twist data.
The construction of the auxiliary realizations adds no condition
on the eigenvalues or automorphy of an arithmetic seed.
% END C69C7
